% BHU PYQ -- Professional Exam Template
\documentclass[11pt,a4paper]{article}

\usepackage[a4paper,top=2.5cm,bottom=2.5cm,left=2.8cm,right=2.8cm,headheight=14pt]{geometry}
\usepackage{amsmath,amssymb}
\usepackage[T1]{fontenc}
\usepackage[utf8]{inputenc}
\usepackage{lmodern}
\usepackage{microtype}
\usepackage{fancyhdr}
\usepackage{enumitem}
\usepackage{hyperref}
\usepackage{lastpage}
\usepackage{parskip}
\usepackage{tikz}
\usetikzlibrary{arrows.meta,shapes,positioning}

\hypersetup{colorlinks=true,urlcolor=black,linkcolor=black,
  pdftitle={B.Sc. (Hons.) Zoology Sem-I ZOB-101 -- BHU PYQ}}

\pagestyle{fancy}
\fancyhf{}
\renewcommand{\headrulewidth}{0.4pt}
\renewcommand{\footrulewidth}{0.4pt}
\fancyhead[L]{\small\textit{Banaras Hindu University}}
\fancyhead[R]{\small\textit{Zoology Paper: ZOB-101}}
\fancyfoot[C]{\small Page \thepage\ of \pageref{LastPage}}

\newlist{parts}{enumerate}{3}
\setlist[parts,1]{label=(\alph*),leftmargin=*,itemsep=4pt,topsep=3pt}
\setlist[parts,2]{label=(\roman*),leftmargin=*,itemsep=2pt,topsep=2pt}
\setlist[parts,3]{label=(\arabic*),leftmargin=*,itemsep=2pt,topsep=2pt}
\newcommand{\pts}[1]{\hfill\textit{[#1]}}

\begin{document}

\begin{center}
  {\large\bfseries BANARAS HINDU UNIVERSITY}\\[2pt]
  {\normalsize B.Sc. (Hons.) Semester I Examination, 2016-17}\\[6pt]
  \rule{0.8\linewidth}{0.5pt}\\[6pt]
  {\large\bfseries Zoology}\\[4pt]
  {\large\bfseries Paper: ZOB-101 --- Animal Diversity \& Animal Form and Function}\\[4pt]
  {\normalsize Time: Three Hours \hfill Full Marks: 35 + 35 = 70}
\end{center}

\rule{\linewidth}{0.4pt}

\smallskip
\noindent\textit{Note: Answer three questions from Section-A and three questions from Section-B. Question No. 1 in both sections is compulsory. Use a separate answer book for each Section. The figures in the right-hand margin indicate marks.}

\bigskip
\begin{center}
  \textbf{SECTION -- A} \\
  \textbf{(Animal Diversity) (Marks: 35)}
\end{center}

\smallskip
\noindent\textit{Note: Answer three questions including Question No. 1 which is compulsory.}

\medskip\noindent\textbf{Question 1.} Answer the following parts:
\begin{parts}
  \item State whether the following statements are True or False. Justify your answer in 2--3 sentences:\pts{$2 \times 3 = 6$}
  \begin{parts}
    \item Balanoglossus belong to class Pterobranchia of Phylum Hemichordata.
    \item Birds are glorified reptiles.
    \item Colloblasts are found in Cnidarians.
  \end{parts}

  \item Distinguish between the following:\pts{$3 \times 1 = 3$}
  \begin{parts}
    \item Protostomes and Deuterostomes
    \item Urodela and Anura
    \item Pseudocoelomate and coelomate.
  \end{parts}

  \item Define the following:\pts{$2 \times 1 = 2$}
  \begin{parts}
    \item Statocyst
    \item Symmetry.
  \end{parts}
\end{parts}

\medskip\noindent\textbf{Question 2.} Answer the following parts:\pts{$7 + 5 = 12$}
\begin{parts}
  \item Classify phylum protozoa up to classes giving characteristic features and examples of each class.
  \item Differentiate between cartilaginous and bony fishes.
\end{parts}

\medskip\noindent\textbf{Question 3.} Answer the following parts:\pts{$7 + 5 = 12$}
\begin{parts}
  \item Classify phylum Mollusca up to classes giving characteristic features and examples of each class.
  \item Describe water vascular system of Echinoderms.
\end{parts}

\medskip\noindent\textbf{Question 4.} Write short notes on any three of the following:\pts{$3 \times 4 = 12$}
\begin{parts}
  \item Skull types in Reptiles
  \item Cetaceans
  \item Cyclostomes
  \item Canal system of Poriferans.
\end{parts}

\pagebreak
\begin{center}
  \textbf{SECTION -- B} \\
  \textbf{(Animal Form and Function) (Marks: 35)}
\end{center}

\smallskip
\noindent\textit{Note: Answer three questions including Question No. 1 which is compulsory.}

\medskip\noindent\textbf{Question 1.} Answer the following parts:
\begin{parts}
  \item State whether the following statements are True or False and give brief reasons in support of your answer:\pts{$2 \times 4 = 8$}
  \begin{parts}
    \item Longer the loop of nephron more water and ions can be reabsorbed.
    \item The presynaptic cell transmits signals from specialized enlarged regions of its axon's branches, called synaptic cleft.
    \item Claspers are accessory sex organ found in birds and mammals.
    \item Motor neurons carry commands from the brain or spinal cord to muscles and glands.
  \end{parts}

  \item Define the following terms:\pts{$1 \times 3 = 3$}
  \begin{parts}
    \item Suspension feeding
    \item Book lungs
    \item Sensory pathways.
  \end{parts}
\end{parts}

\medskip\noindent\textbf{Question 2.} Answer the following parts:\pts{$6 + 6 = 12$}
\begin{parts}
  \item With the help of well labelled diagram explain the structure and functions of mammalian kidney. Describe the role of protonephridia in the process of excretion.
  \item Describe the structure and function of gills in fishes. Explain the role of mammalian lung in gas exchange.
\end{parts}

\medskip\noindent\textbf{Question 3.} Answer the following parts:\pts{$6 + 6 = 12$}
\begin{parts}
  \item Describe the patterns of nervous system in non-chordates with supporting well labelled diagram.
  \item Give an account of the phonoreception in mammals.
\end{parts}

\medskip\noindent\textbf{Question 4.} Write short notes on any three of the following:\pts{$3 \times 4 = 12$}
\begin{parts}
  \item Intracellular digestion
  \item Parthenogenesis
  \item Accessory sex organs in mammals.
\end{parts}

\vfill
\rule{\linewidth}{0.4pt}
\begin{center}\small
\href{https://bhu-exam.vercel.app/}{Click Here}\\[4pt]\href{https://me-aryan.vercel.app/}{Click Here}\\[4pt]\href{https://quashan-me.vercel.app/}{Click Here}
\end{center}

\end{document}
